% =============================================================================
% Section 7 — Gossip Protocol
% =============================================================================
\section{Gossip Protocol: 4-Step Symmetric Synchronisation}
\label{sec:gossip}

\pebble{} uses a \emph{pull-based, symmetric gossip protocol}~\cite{demers1987epidemic,
kermarrec2007gossip} for replication. Every \texttt{gossipIntervalMs} milliseconds,
each node selects one healthy peer (random selection among healthy nodes) and
runs a complete sync session with it.

\subsection{Session Overview}

Each gossip session consists of exactly four steps, alternating message ownership
between the initiator and responder. Both sides make progress in each round.

\begin{figure}[H]
\centering
\begin{tikzpicture}[font=\small, node distance=0.1cm]

  % Column headers
  \node (A) at (1.5, 8.5) {\textbf{Initiator (Node A)}};
  \node (B) at (9.5, 8.5) {\textbf{Responder (Node B)}};
  \draw[thick, gray!40] (3, 8.8) -- (3, -0.3);
  \draw[thick, gray!40] (8, 8.8) -- (8, -0.3);
  \draw[dashed, gray!30] (0.2, 8.2) -- (11, 8.2);

  % Step 1 label
  \node[font=\footnotesize\itshape, left] at (0.2, 7.6) {Step 1};

  % DIGEST A -> B
  \draw[msg arrow] (3, 7.6) -- (8, 7.2)
    node[midway, above, font=\footnotesize]
    {\msg{DIGEST}\{$V_A$\}};
  % DIGEST B -> A
  \draw[msg arrow] (8, 7.2) -- (3, 6.8)
    node[midway, below, font=\footnotesize]
    {\msg{DIGEST}\{$V_B$\}};

  \draw[dashed, gray!30] (0.2, 6.5) -- (11, 6.5);
  \node[font=\footnotesize\itshape, left] at (0.2, 6.1) {Step 2};

  % REQUEST A -> B
  \draw[msg arrow] (3, 6.1) -- (8, 5.7)
    node[midway, above, font=\footnotesize]
    {\msg{REQUEST}\{$\mathrm{missing}(V_B, V_A)$\}};
  % MUTATIONS B -> A
  \draw[msg arrow] (8, 5.7) -- (3, 5.3)
    node[midway, below, font=\footnotesize]
    {\msg{MUTATIONS}\{$m_1, \ldots, m_k$\}};

  \draw[dashed, gray!30] (0.2, 5.0) -- (11, 5.0);
  \node[font=\footnotesize\itshape, left] at (0.2, 4.6) {Step 3};

  % REQUEST B -> A
  \draw[msg arrow] (8, 4.6) -- (3, 4.2)
    node[midway, below, font=\footnotesize]
    {\msg{REQUEST}\{$\mathrm{missing}(V_A, V_B)$\}};
  % MUTATIONS A -> B
  \draw[msg arrow] (3, 4.2) -- (8, 3.8)
    node[midway, above, font=\footnotesize]
    {\msg{MUTATIONS}\{$m_1', \ldots, m_j'$\}};

  \draw[dashed, gray!30] (0.2, 3.5) -- (11, 3.5);
  \node[font=\footnotesize\itshape, left] at (0.2, 3.1) {Step 4};

  % ACK A -> B
  \draw[msg arrow] (3, 3.1) -- (8, 2.7)
    node[midway, above, font=\footnotesize]
    {\msg{ACK}\{\texttt{ok: true}\}};
  % ACK B -> A
  \draw[msg arrow] (8, 2.7) -- (3, 2.3)
    node[midway, below, font=\footnotesize]
    {\msg{ACK}\{\texttt{ok: true}\}};

  % Lifecycle lines
  \draw[thick] (3, 8.2) -- (3, 1.8);
  \draw[thick] (8, 8.2) -- (8, 1.8);

\end{tikzpicture}
\caption{4-step symmetric gossip protocol. After exchanging DIGESTs, each side
requests and receives the mutations it is missing.}
\label{fig:gossip-protocol}
\end{figure}

\subsection{Step-by-Step Description}

\paragraph{Step 1 — DIGEST exchange.}
Both nodes send their current version vectors simultaneously. After this step,
each node holds both $V_A$ and $V_B$ and can independently compute the ranges
that the other is missing.

\paragraph{Step 2 — Initiator requests and receives.}
The initiator computes $\mathrm{missing}(V_B, V_A)$ (what $A$ needs from $B$)
and sends a \msg{REQUEST} message. The responder looks up the requested
mutations in its history and sends them in a \msg{MUTATIONS} message. The
initiator applies each received mutation through the store's
\texttt{applyMutation} pipeline (including signature verification).

\paragraph{Step 3 — Responder requests and receives.}
The responder computes $\mathrm{missing}(V_A, V_B)$ (what $B$ needs from $A$)
and sends a \msg{REQUEST}. The initiator replies with the matching mutations.

\paragraph{Step 4 — ACK.}
Both sides send \msg{ACK} to signal successful completion. The gossip engine
records a successful round and updates the peer's health state.

\subsection{Protocol Properties}

\begin{theorem}[Deadlock freedom]
\label{thm:deadlock-free}
The 4-step gossip protocol is deadlock-free.
\end{theorem}

\begin{proof}
The message sequence is fixed and deterministic: initiator sends, then
responder sends, alternating for exactly 4 sub-steps. Neither side waits for
the other to initiate a send before completing its own step. Each
\texttt{send}/\texttt{receive} call on a connection is backed by a buffered
TCP socket; a send never blocks waiting for the remote side to call \texttt{receive}
on the same message. Since the turn sequence is fixed and finite, no circular
waiting condition can arise. \qed
\end{proof}

\begin{theorem}[Idempotency]
\label{thm:idempotent}
Applying a mutation that has already been applied is a no-op.
\end{theorem}

\begin{proof}
Before applying any mutation $m$, \texttt{applyMutation} checks whether
$\mathrm{id}(m)$ already exists in the history map $H$. If it does, and the
stored content is byte-for-byte identical, the function returns \texttt{false}
without modifying state. If the stored content differs (a sequence collision),
an error is thrown (see Invariant~\ref{inv:uniqueness}). In the normal case,
duplicate mutations are silently dropped, making gossip delivery idempotent. \qed
\end{proof}

\subsection{Gossip Target Selection}

In each round, the \texttt{GossipEngine} calls
\texttt{PeerHealth.selectGossipTarget()}, which returns a random healthy peer
or \texttt{null} if none are available. This provides a simple, unbiased
gossip fanout of 1 per interval. The expected number of rounds for a mutation
to reach all $N$ nodes is $\Oof{N \log N}$ (see \cref{sec:complexity}).

\subsection{Health Monitoring}

The \texttt{Pinger} sends \msg{PING} messages to all peers every
\texttt{pingIntervalMs} milliseconds (default: 2,000\,ms) on the gossip port.
A peer is marked \emph{unhealthy} if no successful communication has been
observed in the past \texttt{unhealthyThresholdMs} milliseconds (default:
10,000\,ms). Unhealthy peers are excluded from gossip target selection but
remain in the configuration; they are automatically restored to healthy status
upon the next successful handshake.
